% Part: turing-machines
% Chapter: undecidability
% Section: representing-tms

\documentclass[../../../include/open-logic-section]{subfiles}

\begin{document}

\olfileid[hi]{tur}{und}{rep}
\olsection{ट्यूरिंग मशीनों का निरूपण}

\begin{explain}
ट्यूरिंग मशीनों और उनके व्यवहार को प्रथम-कोटि तर्कशास्त्र के
!!a{sentence} द्वारा निरूपित करने के लिए हमें उपयुक्त भाषा परिभाषित करनी
होगी। भाषा के दो भाग हैं: मशीन के विन्यासों का वर्णन करने वाले
!!{predicate}, और निष्पादन के चरणों (``क्षणों'') तथा पट्टी के स्थानों को
क्रमांकित करने वाली अभिव्यक्तियाँ।

हम दो प्रकार के !!{predicate} प्रस्तुत करते हैं और दोनों द्विस्थानी हैं:
प्रत्येक अवस्था~$q$ के लिए !!a{predicate}~$\Obj Q_q$, तथा पट्टी के
प्रत्येक प्रतीक~$\sigma$ के लिए !!a{predicate}~$\Obj S_\sigma$। पहले
प्रकार से हम~$M$ की अवस्था और उसके पट्टी-शीर्ष का स्थान वर्णित कर सकते
हैं; दूसरे प्रकार से पट्टी की सामग्री का वर्णन कर सकते हैं।

पट्टी-शीर्ष के स्थानों और निष्पादित चरणों की संख्या को व्यक्त करने के
लिए हमें संख्याएँ व्यक्त करने की कोई रीति चाहिए। इसके लिए !!a{constant}
~$\Obj 0$ और $1$-स्थानी उत्तरवर्ती फलन~$\prime$ का उपयोग होता है। परिपाटी
के अनुसार इसे अपने आगत के \emph{बाद} लिखा जाता है (और हम कोष्ठक छोड़
देते हैं)।

प्रत्येक संख्या $n$ के लिए एक विहित पद~$\num{n}$ है, जो~$n$ का
\emph{अंक-पद} है और~$\Lang L_M$ में उसे निरूपित करता है। $\num{0}$,
$\Obj 0$ है; $\num{1}$, $\Obj 0'$ है; $\num{2}$, $\Obj 0''$ है; और इसी
प्रकार आगे। अधिक औपचारिक रूप से:
\begin{align*}
\num{0} & = \Obj 0 \\
\num{n+1} &= \num{n}'
\end{align*}
पद $\num{0}$, अर्थात् $\Obj 0$, पट्टी के सबसे बाएँ स्थान और प्रथम
निष्पादन-चरण से पहले के समय (आरंभिक विन्यास) दोनों को नाम देता है। पद
$\num{1}$, अर्थात् $\Obj 0'$, सबसे बाएँ खाने की दाईं ओर के खाने और प्रथम
निष्पादन-चरण के बाद के समय को नाम देता है; और इसी प्रकार आगे।

हम पट्टी के स्थानों के क्रम (जहाँ इसका अर्थ ``के बाईं ओर'' है) और
निष्पादन-चरणों के क्रम (जहाँ इसका अर्थ ``से पहले'' है) दोनों को व्यक्त
करने के लिए !!a{predicate}~$<$ भी प्रस्तुत करते हैं।

भाषा निर्धारित हो जाने पर हम $!T(M, w)$ के ``अभिगृहीत'' सूचीबद्ध करते
हैं, अर्थात् वे !!{sentence} जो साथ मिलकर आगत~$w$ पर चलने वाली~$M$ का
व्यवहार वर्णित करते हैं। कुछ !!{sentence}, $\Obj 0$, $\prime$, और $<$
पर शर्तें लगाएँगे; कुछ !!{sentence} आगत-विन्यास का वर्णन करेंगे; और कुछ
!!{sentence} यह बताएँगे कि किसी विशेष निर्देश को निष्पादित करने के बाद
$M$ का विन्यास क्या है।
\end{explain}

\begin{defn}
  \ollabel{defn:tm-descr}
दी गई ट्यूरिंग मशीन $M = \tuple{Q, \Sigma, q_0, \delta}$ के लिए
भाषा~$\Lang L_M$ में ये होते हैं:
\begin{enumerate}
\item प्रत्येक अवस्था~$q \in Q$ के लिए द्विस्थानी !!{predicate}
  $\Obj Q_q(x, y)$। सहज रूप से, $\Obj Q_q(\num{m}, \num{n})$ यह व्यक्त
  करता है कि ``$n$ चरणों के बाद $M$, अवस्था~$q$ में है और $m$वाँ खाना
  पढ़ रही है।''
\item प्रत्येक प्रतीक~$\sigma\in \Sigma$ के लिए द्विस्थानी !!{predicate}
  $\Obj S_\sigma(x, y)$। सहज रूप से, $\Obj S_\sigma(\num{m}, \num{n})$
  यह व्यक्त करता है कि ``$n$ चरणों के बाद $m$वें खाने में प्रतीक
  ~$\sigma$ है।''
\item एक !!{constant} $\Obj 0$
\item एक-स्थानी !!{function} $\prime$
\item द्विस्थानी !!{predicate} $<$
\end{enumerate}
\end{defn}

आगत $w = \sigma_{i_1}\dots\sigma_{i_k}$ पर ट्यूरिंग मशीन~$M$ की क्रिया
का वर्णन करने वाले !!{sentence} निम्न हैं:
\begin{enumerate}
\item संख्याओं और~$<$ का वर्णन करने वाले अभिगृहीत:
\begin{enumerate}
\item यह कहने वाला !!^a{sentence} कि प्रत्येक संख्या अपने उत्तरवर्ती से
छोटी है:
\[
\lforall[x][x < x']
\]
\item यह सुनिश्चित करने वाला !!^a{sentence} कि $<$ संक्रामक है:
\[
\lforall[x][\lforall[y][\lforall[z][
      ((x < y \land y < z) \lif x < z)]]]
\]
\end{enumerate}
\item आगत-विन्यास का वर्णन करने वाले अभिगृहीत:
\begin{enumerate}
\item $0$ चरणों के बाद---मशीन के आरंभ होने से पहले---$M$ आरंभिक
  अवस्था~$q_0$ में है और खाना~$1$ पढ़ रही है:
\[
\Obj Q_{q_0}(\num{1}, \num{0})
\]
\item पहले $k+1$ खानों में प्रतीक $\TMendtape$, $\sigma_{i_1}$, \dots,
  $\sigma_{i_k}$ हैं:
\[
\Obj S_\TMendtape(\num{0}, \num{0}) \land
\Obj S_{\sigma_{i_1}}(\num{1}, \num{0}) \land
\dots \land
\Obj S_{\sigma_{i_k}}(\num{k}, \num{0})
\]
\item अन्यथा पट्टी रिक्त है:
\[
\lforall[x][(\num{k} < x \lif \Obj S_\TMblank(x, \num{0}))]
\]
\end{enumerate}
\item एक विन्यास से अगले विन्यास तक संक्रमण का वर्णन करने वाले अभिगृहीत:

निम्न के लिए $!A(x, y)$ को इस रूप वाले सभी !!{sentence} का संयोजन मानें:
\[
\lforall[z][
  (((z < x \lor x < z) \land \Obj S_\sigma(z, y))
  \lif \Obj S_\sigma(z, y'))]
\]
जहाँ $\sigma \in \Sigma$। हम $!A(\num{m},\num{n})$ का उपयोग यह व्यक्त
करने के लिए करते हैं कि ``खाना~$m$ को छोड़कर, $n+1$ चरणों के बाद पट्टी
वैसी ही है जैसी $n$ चरणों के बाद थी।''
\begin{enumerate}
\item \ollabel{rep-right} प्रत्येक निर्देश $\delta(q_i, \sigma) =
  \tuple{q_j, \sigma', \TMright}$ के लिए निम्न !!{sentence}:
\begin{align*}
& \lforall[x][\lforall[y][(
   (\Obj Q_{q_i}(x, y) \land \Obj S_{\sigma}(x, y)) \lif {}]] \\
&\qquad   (\Obj Q_{q_j}(x', y') \land \Obj S_{\sigma'}(x, y') \land
!A(x, y)))
\end{align*}
यह कहता है कि यदि~$y$ चरणों के बाद मशीन अवस्था~$q_i$ में है और
प्रतीक~$\sigma$ रखने वाला खाना~$x$ पढ़ रही है, तो $y+1$ चरणों के बाद वह
खाना~$x+1$ पढ़ रही है, अवस्था~$q_j$ में है, खाना~$x$ में अब~$\sigma'$
है, और~$x$ के अतिरिक्त प्रत्येक खाने में वही प्रतीक है जो उसमें~$y$
चरणों के बाद था।

\item \ollabel{rep-left} प्रत्येक निर्देश $\delta(q_i, \sigma) =
  \tuple{q_j, \sigma', \TMleft}$, के लिए निम्न !!{sentence}:
\begin{align*}
& \lforall[x][\lforall[y][
    ((\Obj Q_{q_i}(x', y) \land \Obj S_{\sigma}(x', y)) \lif {}]]\\
& \qquad   (\Obj Q_{q_j}(x, y') \land \Obj S_{\sigma'}(x', y') \land
!A(x, y))) \land {}\\
& \lforall[y][((\Obj Q_{q_i}(\num{0}, y) \land \Obj S_{\sigma}(\num{0},
    y)) \lif {}]\\
& \qquad (\Obj Q_{q_j}(\num{0}, y') \land \Obj S_{\sigma'}(\num{0},
  y') \land !A(\num{0}, y)))
\end{align*}
यह कैसे काम करता है, इस पर कुछ क्षण विचार करें: अब हम ``यदि खाना~$x$
पढ़ रही है \dots'' से नहीं, बल्कि ``यदि खाना $x+1$ पढ़ रही है \dots''
से आरंभ करते हैं। बाईं ओर जाने का अर्थ है कि अगले चरण में मशीन खाना~$x$
पढ़ रही होगी। किंतु जिस खाने पर लिखा जाता है वह~$x+1$ है। हम ऐसा इसलिए
करते हैं क्योंकि हमारे पास घटाव अथवा पूर्ववर्ती फलन नहीं है।

ध्यान दें कि $x+1$ रूप वाली संख्याएँ $1$, $2$, \dots हैं; अर्थात् इसमें
वह स्थिति सम्मिलित नहीं है जिसमें मशीन खाना~$0$ पढ़ रही है और उसे बाईं
ओर जाना है (जो निश्चय ही वह नहीं कर सकती---वह वहीं रहती है)। इस विशेष
स्थिति को दूसरा संयोजन समाहित करता है: वह कहता है कि यदि $y$ चरणों के
बाद मशीन अवस्था~$q_i$ में खाना~$0$ पढ़ रही है और खाना~$0$ में
प्रतीक~$\sigma$ है, तो $y+1$ चरणों के बाद भी वह खाना~$0$ पढ़ रही है,
अब अवस्था~$q_j$ में है, खाना~$0$ पर प्रतीक $\sigma'$ है, और खाना~$0$ के
अतिरिक्त खानों में वही प्रतीक हैं जो उनमें $y$ चरणों के बाद थे।
\item \ollabel{rep-stay} प्रत्येक निर्देश $\delta(q_i, \sigma) =
  \tuple{q_j, \sigma', \TMstay}$ के लिए निम्न !!{sentence}:
\begin{align*}
& \lforall[x][\lforall[y][(
   (\Obj Q_{q_i}(x, y) \land \Obj S_{\sigma}(x, y)) \lif {}]] \\
&\qquad   (\Obj Q_{q_j}(x, y') \land \Obj S_{\sigma'}(x, y') \land
!A(x, y)))
\end{align*}
\end{enumerate}
\end{enumerate}
ट्यूरिंग मशीन~$M$ और आगत~$w$ के लिए ऊपर के सभी !!{sentence} के संयोजन
को $!T(M, w)$ मानें।

यह व्यक्त करने के लिए कि~$M$ अंततः रुकती है, हमें ऐसा !!{sentence} खोजना
होगा जो कहे, ``कुछ चरणों के बाद संक्रमण फलन अपरिभाषित होगा।'' $X$ को
उन सभी युग्मों $\tuple{q, \sigma}$ का समुच्चय मानें जिनके लिए
~$\delta(q, \sigma)$ अपरिभाषित है। तब $!E(M, w)$ को निम्न !!{sentence}
मानें:
\[
\lexists[x][\lexists[y][(\bigvee_{\tuple{q, \sigma} \in
      X}(\Obj Q_q(x, y) \land \Obj S_\sigma(x, y)))]]
\]

यदि हम निर्दिष्ट विराम-अवस्था~$h$ वाली ट्यूरिंग मशीन का उपयोग करें, तो
यह और भी सरल है: तब निम्न !!{sentence}~$!E(M, w)$
\[
\lexists[x][\lexists[y][\Obj Q_h(x, y)]]
\]
व्यक्त करता है कि मशीन अंततः रुकती है।

\begin{prop}
\ollabel{prop:mlessk}
यदि $m < k$, तो $!T(M, w) \Entails \num{m} < \num{k}$।
\end{prop}

\begin{proof}
अभ्यास।
\end{proof}

\begin{prob}
\olref[tur][und][rep]{prop:mlessk} सिद्ध करें।
(संकेत: $k-m$ पर आगमन का उपयोग करें।)
\end{prob}

\end{document}
